% --------------------------------------------------------------------

\section{Name-based identity vs content-addressing}

This section describes the problem and the proposed solution at a
  high level.  We first show how Lean's name-based identity couples
  references to project structure, file layout, and naming conventions,
  making them fragile under refactoring and opaque across project
  boundaries.  We then introduce content-addressing---hashing the
  canonicalized type expression rather than the name---and state its
  key properties.

\paragraph{Name-based Identity}
% =====================================================================

To use a theorem in Lean~4, a programmer writes something like:

\begin{lstlisting}[caption={Using a theorem by its fully qualified name}]
import Mathlib.Data.Finset.Sum

example : forall (s : Finset Nat) (f g : Nat -> Int),
    s.sum (f + g) = s.sum f + s.sum g :=
  Finset.sum_add_distrib
\end{lstlisting}

This single reference, \texttt{Finset.sum\_add\_distrib}, encodes three
pieces of information beyond the mathematics: (1)~the \textbf{project}
(Mathlib), (2)~the \textbf{file}
(\texttt{Mathlib/Data/Finset/Sum.lean}, needed for the \texttt{import}),
and (3)~the \textbf{namespace path}
(\texttt{Finset.sum\_add\_distrib}, a human-chosen convention).
None of these are intrinsic to the mathematical statement.  The theorem
says that finite sums distribute over addition---a fact that is true
independently of what we call it or where we store it.

When any of these coordinates change, references break.  A theorem
renamed for consistency (\texttt{sum\_add\_distrib} $\to$
\texttt{sum\_add}), a declaration relocated to a different module
(\texttt{Data.Finset.Sum} $\to$ \texttt{Algebra.BigOps}), a namespace
hierarchy flattened or reorganized, a project forked with new
prefixes---each is a no-op mathematically, but a breaking change
syntactically.  The cost is borne by every downstream project: CI
pipelines fail, students' homework stops compiling, automated tooling
must track name migrations.

Within a single library like Mathlib, naming conventions and import
paths are at least centrally managed.  Across independent projects, the
situation is worse.  There is no standardized way for project~A to
refer to a result proved in project~B, other than adding a direct
dependency on B's entire package---including its specific version, file
layout, and naming scheme.  If two independent projects prove the same
theorem under different names, there is no mechanism to recognize that
they state the same mathematical fact.  Each project is an island, and
the only bridges are explicit package dependencies with name-level
coupling.

This problem will only grow.  Mathlib is evolving into the shared
foundation that virtually all Lean projects depend on, but the
ecosystem is expanding well beyond pure mathematics: formalization
projects in cryptography, blockchain protocol verification, physics,
and other applied domains are emerging, each building on Mathlib and
on each other.  As the number of independent projects grows, so does
the cost of name-level coupling---every cross-project reference is a
fragile link that any upstream refactoring can break.


% --------------------------------------------------------------------
\paragraph{Content-Addressing}
% --------------------------------------------------------------------

The central idea is simple: \emph{hash the mathematics, not the name.}

Given a declaration with a type expression (its mathematical statement),
we: (1)~\textbf{canonicalize} the type expression to remove
representation artifacts, (2)~\textbf{serialize} the canonical
expression to a byte sequence, and (3)~\textbf{hash} the byte sequence
with SHA-256.  The resulting 256-bit hash is the declaration's
\emph{content address}.
For theorems, the content address is therefore determined entirely by
the statement---the proof term is never hashed.
For propositions, this is inevitable given proof irrelevance.

Content addresses have four key properties.  They are \emph{stable}:
the address does not change when a declaration is renamed, moved, or
reorganized.  They are \emph{universal}: no project, file, or import
context is needed, and no coordination or central authority is required.
They are \emph{self-authenticating}: anyone can verify an address by
recomputing the hash from the declaration's type.  And they are
\emph{collision-resistant}: two mathematically different statements will
(with overwhelming probability) have different addresses.

An important caveat: content addresses identify declarations that share
the same \emph{canonical type}, but this does not make them freely
interchangeable.  Declarations carry metadata (attributes, instance
registrations) that the hash does not capture, and definitions with the
same type may compute differently.  We discuss these limits in
Section~\ref{sec:equiv}.

\subsection{Features}

  \paragraph{Search by statement, not by name}
  \label{rem:search}
  Content-addressing inverts the usual discovery workflow.  To determine
  whether a result is already known---proved, published as an open
  conjecture, or simply mentioned in a registry---one does not need to
  guess the name a library author might have chosen, nor search through
  documentation by keyword.  Instead, it suffices to \emph{write the
  statement} as a Lean type expression, compute its content address
  (a single call to the hashing pipeline), and look up that address
  in the registry.  A match means the exact mathematical fact is already
  recorded, regardless of the name, file, or project it lives in.
  A miss means no project in the registry has published that statement
  at the chosen canonicalization level.  The search is determined
  entirely by the mathematics---no text matching, no naming conventions,
  no knowledge of library structure required.

  \paragraph{Address is stable under proof improvement}
  [Stable under proof improvement]
  \label{rem:proof-improvement}
  Because the content address depends only on the statement, improving a
  proof---making it shorter, faster, or more elegant---does not change
  the address.  The mathematical ``API'' is stable even as
  implementations improve.  A library can silently replace a slow proof
  with a faster one and no downstream content address is affected.
  It's the analogue of "you can refactor the implementation without changing the interface" in software.

  \paragraph{Content addresses as permanent citations}
  \label{rem:citations}
  Academic papers about formal mathematics currently cite theorems by
  library name (e.g., \texttt{Finset.sum\_add\_distrib}).  These
  citations break silently when libraries refactor.  A content address
  serves as a permanent reference to a mathematical fact.
  A paper citing \texttt{ca:a1b2c3\ldots} remains
  valid regardless of future renaming or reorganization.

  \paragraph{Memoization for proof search (AI angle)}
  \label{rem:memoization}
  In automated or AI-assisted proof search, the same goal type often
  arises in different proof contexts.  A content address canonically
  identifies the goal independently of the surrounding proof state.
  This enables a global lookup table: before attempting synthesis, check
  whether the goal's content address already appears in the registry.
  If it does, the proof can be retrieved rather than re-discovered.

  \paragraph{Version-independent references}
  \label{rem:versions}
  A content address for a theorem is the same whether computed from
  Mathlib v4.3 or v4.9, provided the mathematical statement has not
  changed.  This decouples mathematical identity from library
  versioning: a project can depend on a \emph{mathematical fact} rather
  than on a specific version of a specific package.

  \paragraph{Progress tracking for formalization campaigns}
  \label{rem:campaigns}
  A formalization effort (e.g., all theorems in a textbook chapter) can
  be specified as a list of content addresses registered as open points.
  Progress is then the fraction of addresses whose status has moved from
  \texttt{open} to \texttt{proved}, a well-defined, cross-project
  metric that updates automatically via the status cascade.

  \paragraph{Toward granular building.}
  Lean's build system (Lake) operates at the granularity of entire
  modules: importing a single declaration requires compiling the full
  module and all its transitive dependencies.  This monolithic model
  compounds the cost of name-based coupling---a project that needs one
  lemma from a large module pays for thousands of unrelated declarations.
  Content-addressing opens a path toward finer granularity.  If
  declarations are individually addressable by content, a build system
  could in principle fetch and verify only the specific results a project
  actually uses, rather than pulling in entire module trees.  Granular
  addressing would lead naturally to granular imports: referencing
  individual facts rather than files.  This is a harder problem in a
  compiled ecosystem like Lean, where elaboration, compilation, and
  kernel-checking are tightly coupled at the module level, but
  content-addressing provides the necessary first step---a stable,
  module-independent name for each declaration.
  
% =====================================================================
